\documentclass{article}
\usepackage{amsmath}
\newcommand{\R}{\mathbb{R}}
\newcommand{\norm}[1]{\left\lVert #1 \right\rVert}
\newtheorem{theorem}{Theorem}
\newtheorem{lemma}{Lemma}
\title{Every construct, in \LaTeX}
\author{Ann Writer \and Bo Reader}
\date{November 2026}

\begin{document}
\maketitle

\begin{abstract}
A short abstract with \emph{emphasis} and a citation \cite{knuth84}.
\end{abstract}

\section{Introduction}\label{sec:intro}
Plain text with \textbf{bold}, \textit{italic}, \emph{emphasis}, \underline{underlined},
\texttt{typewriter}, \textsc{small caps}, \sout{struck}, and \textbf{bold with \emph{emphasis} inside}.
A tie~stays a space, a comment % this part is a comment
ends the line, and an escaped \% sign, \& and \$ stay text.
A link \href{https://example.com}{to an example} and a bare \url{https://example.org}.
A footnote\footnote{The footnote's words.} and a line break \\ in a paragraph.

Inline maths $x \in \R$, \(\norm{v} = 1\), and display maths:
\[
  \sum_{i=1}^{n} i = \frac{n(n+1)}{2}
\]

\subsection{Numbered things}\label{sec:numbered}
\begin{equation}\label{eq:first}
  a^2 + b^2 = c^2
\end{equation}
As Section~\ref{sec:intro} and equation~\eqref{eq:first} say, and \autoref{eq:second} too.
\begin{align}
  x &= 1 \label{eq:second} \\
  y &= 2 \nonumber \\
  z &= 3
\end{align}
\begin{equation*}
  \text{unnumbered}
\end{equation*}

\subsubsection{Lists}
\begin{itemize}
  \item First \textbf{item}
  \item Second, with a nested list:
  \begin{enumerate}
    \item One
    \item Two \label{item:two}
  \end{enumerate}
  \item[$\star$] A custom marker
\end{itemize}
\begin{description}
  \item[Term] Its description.
\end{description}

\begin{theorem}[Pythagoras]\label{thm:p}
For a right triangle, $a^2 + b^2 = c^2$.
\end{theorem}
\begin{proof}
By Theorem~\ref{thm:p} itself, and Lemma~\ref{lem:x}.
\end{proof}
\begin{lemma}\label{lem:x}
A lemma.
\end{lemma}

\begin{figure}[h]
  \centering
  \includegraphics[width=0.5\textwidth]{pictures/plot.pdf}
  \caption{A plot, as Figure~\ref{fig:plot} says.}\label{fig:plot}
\end{figure}

\begin{table}[h]
  \centering
  \begin{tabular}{l|c|r}
    Left & Centre & Right \\
    \hline
    a & \textbf{b} & $c$ \\
    \multicolumn{2}{c}{spanned} & d \\
  \end{tabular}
  \caption{A table.}\label{tab:t}
\end{table}

\begin{verbatim}
verbatim \textbf{stays} as written
\end{verbatim}

\paragraph{A paragraph heading} with words after it, citing \cite{lamport94,knuth84} and \citep{knuth84}.

\begin{thebibliography}{9}
\bibitem{lamport94} Leslie Lamport. \emph{LaTeX}. 1994.
\bibitem{knuth84} Donald Knuth. \emph{The TeXbook}. 1984.
\end{thebibliography}
\end{document}
